\documentclass{beamer}

\usetheme{Berlin}
\usecolortheme{Orchid}
\useoutertheme{miniframes}
\setbeamertemplate{navigation symbols}{}

\usepackage[utf8]{inputenc}
\usepackage[T1]{fontenc}
\usepackage{amsmath}
\usepackage{amssymb}
\usepackage{booktabs}
\usepackage{graphicx}
\graphicspath{{../images/}}

\usepackage{xcolor}
\usepackage{listings}
\usepackage{hyperref}
\usepackage{tikz}
\usetikzlibrary{arrows.meta,positioning,shapes.geometric,calc}

\lstset{
  language=Java,
  basicstyle=\ttfamily\small,
  keywordstyle=\color{blue},
  commentstyle=\color{gray},
  stringstyle=\color{teal},
  showstringspaces=false,
  columns=fullflexible,
  breaklines=true
}

% Per-slide source line (references shown on the slide itself).
\newcommand{\slidesources}[1]{%
  \vfill
  {\tiny\color{gray}\raggedright\sloppy\parbox{\linewidth}{Sources: #1}\par}%
}

% Unit-wise source strings for this subject (used by slides/notes).
% Path is relative to the lecture `latex/` folder (the compilation CWD).
% OOPS with Java (CSEG1044) — unit-wise sources
%
% These are short, student-facing reference strings intended to be used with:
%   - \slidesources{...}  (in beamer slides)
%   - \notesources{...}   (in lecture notes)
%
% Style inspiration: other subjects in My-Subjects/ use semicolon-separated sources
% and include an "accessed" date for web documentation.

\newcommand{\OOPSUnitISources}{Schildt, \textit{Java: A Beginner's Guide} (10e), McGraw-Hill (Java fundamentals + OOP basics).; Oracle Java Tutorials: Getting Started + Language Basics + Classes and Objects (accessed \today).; Java SE API docs (e.g., \texttt{String}, \texttt{Scanner}) (accessed \today).}

\newcommand{\OOPSUnitIISources}{Schildt, \textit{Java: The Complete Reference} (12e), McGraw-Hill (classes, inheritance, packages, interfaces).; Bloch, \textit{Effective Java} (3e), Addison-Wesley, 2018 (design choices; interfaces vs abstract classes).; Oracle Java Tutorials: Classes and Objects + Interfaces + Packages (accessed \today).; Java SE API docs (e.g., \texttt{Comparable}, \texttt{Runnable}) (accessed \today).}

\newcommand{\OOPSUnitIIISources}{Schildt, \textit{Java: The Complete Reference} (12e) (nested/inner classes, exceptions, threads, I/O).; Oracle Java Tutorials: Nested Classes + Exceptions + Concurrency + Basic I/O (accessed \today).; Java SE API docs (e.g., \texttt{Thread}, \texttt{AutoCloseable}) (accessed \today).}

\newcommand{\OOPSUnitIVSources}{Schildt, \textit{Java: The Complete Reference} (12e) (generics, lambdas, Swing, JDBC).; Bloch, \textit{Effective Java} (3e) (generics + lambdas).; Oracle Java Tutorials: Generics + Lambda Expressions + Swing + JDBC Basics (accessed \today).; Java SE API docs (e.g., \texttt{java.util.function}, \texttt{javax.swing}, \texttt{java.sql}) (accessed \today).}

\newcommand{\OOPSUnitVSources}{Schildt, \textit{Java: The Complete Reference} (12e) (Collections Framework + wrapper classes).; Oracle Java docs: Collections Framework overview (accessed \today).; Java SE API docs (\texttt{java.util} collections; wrapper classes in \texttt{java.lang}) (accessed \today).}

\newcommand{\OOPSUnitVISources}{Git SCM: \textit{Pro Git} (2e) (version control workflow), accessed \today.; GitHub Docs (repositories, issues, pull requests), accessed \today.; Oracle Java Tutorials: Swing + JDBC (accessed \today).; JUnit 5 User Guide (accessed \today).; Apache Maven Project: Getting Started (accessed \today).; Gradle User Manual: Getting Started (accessed \today).; UML class diagrams: UML 2.x overview/spec (accessed \today).}

\newcommand{\OOPSMiscWhyInterfacesSources}{Schildt, \textit{Java: The Complete Reference} (12e) (interfaces).; Bloch, \textit{Effective Java} (3e) (prefer interfaces where appropriate).; Oracle Java Tutorials: Interfaces (accessed \today).; Java SE API docs (e.g., \texttt{Comparable}, \texttt{Runnable}) (accessed \today).}



\title[OOPS with Java]{Object Oriented Programming (CSEG1044)}
\subtitle{Unit 05 -- Lecture 01: Collections Framework Overview + Iteration + List}
\begin{document}

\begin{frame}[fragile]
  \titlepage
  \vspace{0.5em}
  \slidesources{\OOPSUnitVSources}
\end{frame}

\begin{frame}{Quick Links}
  \centering
  \hyperlink{sec:overview}{\beamerbutton{Overview}}\hspace{0.6em}
  \hyperlink{sec:concepts}{\beamerbutton{Key Topics}}\hspace{0.6em}
  \hyperlink{sec:demo}{\beamerbutton{Demo}}\hspace{0.6em}
  \hyperlink{sec:summary}{\beamerbutton{Summary}}
  \slidesources{\OOPSUnitVSources}
\end{frame}

\begin{frame}{Agenda}
  \tableofcontents
  \slidesources{\OOPSUnitVSources}
\end{frame}

\section{Overview}
\label{sec:overview}

\begin{frame}{Lecture Snapshot}
\begin{itemize}[<+->]
  \item \textbf{Unit focus:} Collections Framework and Wrapper Classes
  \item \textbf{Course thread:} Use Collections Framework Overview + Iteration + List to manage in-memory project data more cleanly and predictably inside Campus Resource Manager.
\end{itemize}
  \slidesources{\OOPSUnitVSources}
\end{frame}

\begin{frame}{Recall Before You Start}
\begin{itemize}[<+->]
  \item \textbf{Generics}: Collections are safer and clearer once typed containers are already familiar. Main reference: \texttt{Unit 04 / Lecture 01 slides/notes}.
  \item \textbf{Comparison logic}: Sorting becomes easier once small comparison behavior is familiar. Main reference: \texttt{Unit 04 / Lecture 03 slides/notes}.
\end{itemize}
  \slidesources{\OOPSUnitVSources}
\end{frame}


\begin{frame}{Learning Outcomes}
  \begin{itemize}[<+->]
    \item Explain \texttt{Iterable}, \texttt{Collection}, and \texttt{Iterator}.
    \item Choose \texttt{ArrayList} vs \texttt{LinkedList} for common use cases.
    \item Sort a list using \texttt{Comparable} and \texttt{Comparator} (with lambdas).
  \end{itemize}
  \slidesources{\OOPSUnitVSources}
\end{frame}

\begin{frame}{Key Terms to Know}
\small
\begin{columns}[t]
\column{0.48\textwidth}
\begin{itemize}
  \item \textbf{Collection}: group of elements (root interface for many collections).
  \item \textbf{Iterable}: can be used in for-each loop (provides \texttt{iterator()}).
  \item \textbf{Iterator}: cursor used to traverse a collection safely.
\end{itemize}
\column{0.48\textwidth}
\begin{itemize}
  \item \textbf{List}: ordered collection with duplicates allowed and index-based access.
  \item \textbf{Comparable}: natural ordering (\texttt{compareTo}).
  \item \textbf{Comparator}: custom ordering (\texttt{compare}).
\end{itemize}
\end{columns}
  \slidesources{\OOPSUnitVSources}
\end{frame}

\begin{frame}{Study Flow for This Lecture}
\begin{enumerate}[<+->]
  \item Refresh the prerequisite bridge before reading new ideas in isolation.
  \item Study the main build in this order: \textit{Core interfaces: \texttt{Iterable}, \texttt{Collection}, \texttt{Iterator}} $\rightarrow$ \textit{List: ArrayList vs LinkedList} $\rightarrow$ \textit{Sorting: Comparable vs Comparator}.
  \item Trace the worked example and connect each line back to one concept frame.
  \item Attempt the mini exercises before moving to the recap and exit check.
\end{enumerate}
  \slidesources{\OOPSUnitVSources}
\end{frame}


\section{Key Topics}
\label{sec:concepts}

\begin{frame}{Unit 05: Collections and Wrapper Class}
  \begin{itemize}[<+->]
    \item Lecture focus: Collections Framework Overview + Iteration + List
  \end{itemize}
  \slidesources{\OOPSUnitVSources}
\end{frame}

\begin{frame}{Today: Roadmap}
  \begin{itemize}[<+->]
    \item Core interfaces: `Collection`, `Iterable`, `Iterator`
    \item List interface; `ArrayList` vs `LinkedList` (when to use which)
    \item Sorting with `Comparable`/`Comparator` (use lambdas where helpful)
  \end{itemize}
  \slidesources{\OOPSUnitVSources}
\end{frame}

\begin{frame}{Collections Framework (why)}
  \begin{itemize}[<+->]
    \item Standard data structures + algorithms (reduce bugs).
    \item Works well with generics: \texttt{List<Student>}.
    \item Includes sorting utilities and iterators.
  \end{itemize}
  \slidesources{\OOPSUnitVSources}
\end{frame}

\begin{frame}{Iteration (safe removal)}
  \begin{itemize}[<+->]
    \item for-each loop uses an \texttt{Iterator} internally.
    \item Removing inside for-each can throw \texttt{ConcurrentModificationException}.
    \item Use \texttt{Iterator.remove()} when you need to delete while looping.
  \end{itemize}
  \slidesources{\OOPSUnitVSources}
\end{frame}

\begin{frame}{ArrayList vs LinkedList (quick)}
  \begin{itemize}[<+->]
    \item \texttt{ArrayList}: fast random access; best default choice.
    \item \texttt{LinkedList}: slower random access; useful as queue/deque.
  \end{itemize}
  \slidesources{\OOPSUnitVSources}
\end{frame}

\begin{frame}{Sorting: Comparable vs Comparator}
  \begin{itemize}[<+->]
    \item \texttt{Comparable}: natural ordering (\texttt{compareTo}) in the class.
    \item \texttt{Comparator}: custom ordering (\texttt{compare}) outside the class.
    \item Lambdas make comparators short and readable.
  \end{itemize}
  \slidesources{\OOPSUnitVSources}
\end{frame}

\begin{frame}{Checkpoint and Reflection}
\begin{block}{Reflection prompt}
Where would \textit{Collections Framework Overview + Iteration + List} matter most in Campus Resource Manager, and which design choice would you justify using this lecture's ideas?
\end{block}
\begin{block}{Quick self-check}
Which part needs one more example before you move on: \textit{Core interfaces: \texttt{Iterable}, \texttt{Collection}, \texttt{Iterator}}, \textit{List: ArrayList vs LinkedList}, the worked example, or the recap?
\end{block}
  \slidesources{\OOPSUnitVSources}
\end{frame}

\section{Demo / Example}
\label{sec:demo}

\begin{frame}[fragile,shrink=8]{Example: List + iterator + sorting (1/2)}
\begin{lstlisting}
import java.util.*;

class Student {
    final String rollNo;
    final double cgpa;
    Student(String rollNo, double cgpa) { this.rollNo = rollNo; this.cgpa = cgpa; }
    double getCgpa() { return cgpa; }
    public String toString() { return rollNo + ":" + cgpa; }
}

public class Demo {
    public static void main(String[] args) {
        List<Student> list = new ArrayList<>();
        list.add(new Student("21CS1002", 9.1));
        list.add(new Student("21CS1001", 8.2));
        list.add(new Student("21CS1003", 3.9));
\end{lstlisting}
  \slidesources{\OOPSUnitVSources}
\end{frame}

\begin{frame}[fragile]{Example: List + iterator + sorting (2/2)}
\begin{lstlisting}
// continued...
        Iterator<Student> it = list.iterator();
        while (it.hasNext()) {
            if (it.next().getCgpa() < 4.0) it.remove();
        }

        list.sort((a, b) -> Double.compare(b.getCgpa(), a.getCgpa())); // desc
        System.out.println(list);
    }
}
\end{lstlisting}
  \slidesources{\OOPSUnitVSources}
\end{frame}

\begin{frame}{Mini Exercises (2--3 mins)}
  \begin{enumerate}[<+->]
    \item Why use \texttt{Iterator.remove()} when deleting during traversal?
    \item Which list is best for random access: \texttt{ArrayList} or \texttt{LinkedList}?
    \item Difference between \texttt{Comparable} and \texttt{Comparator}?
  \end{enumerate}
  \slidesources{\OOPSUnitVSources}
\end{frame}

\begin{frame}{Watch-outs}
\begin{itemize}[<+->]
  \item Ignoring generics and using raw collections.
  \item Changing a collection unsafely while iterating.
  \item Choosing a List when the real requirement is uniqueness.
\end{itemize}
  \slidesources{\OOPSUnitVSources}
\end{frame}

\section{Summary}
\label{sec:summary}

\begin{frame}{Key Takeaways}
  \begin{itemize}[<+->]
    \item Lists store ordered data and allow duplicates.
    \item Iterators enable safe traversal and removal.
    \item Sorting is done using \texttt{Comparable} or \texttt{Comparator}.
  \end{itemize}
  \slidesources{\OOPSUnitVSources}
\end{frame}

\begin{frame}{Checkpoint Questions}
  \begin{enumerate}[<+->]
    \item What is the purpose of \texttt{Iterator}?
    \item When do you choose \texttt{LinkedList} over \texttt{ArrayList}?
    \item How do you sort a list by a custom rule?
  \end{enumerate}
  \slidesources{\OOPSUnitVSources}
\end{frame}

\begin{frame}{Next Study Steps}
\begin{itemize}[<+->]
  \item Revisit the recall references if any older idea still feels weak.
  \item Rework the lecture example without looking at the notes first.
  \item Attempt the self-study practice and then answer the exit questions in full sentences.
  \item Apply the idea once inside Campus Resource Manager to make the concept stick.
\end{itemize}
  \slidesources{\OOPSUnitVSources}
\end{frame}

\begin{frame}{Next Lecture Preview}
  \textbf{Next:} Set/SortedSet + Map/SortedMap
  \vspace{0.6em}
  \begin{itemize}[<+->]
    \item Uniqueness (Set) and key-value storage (Map) with hashing basics.
  \end{itemize}
  \slidesources{\OOPSUnitVSources}
\end{frame}

\end{document}
